% V2 PDF物理页118原书思考题16.1--16.9；图16.18理想电表条件待专家核验。
\section*{第16章\quad 恒定电流：思考题解析}
\paragraph{思考题 16.1}
导体可有局部电场但处于开路或电荷无净漂移，故总电流为零；若稳态、均匀欧姆导体 \(\boldsymbol J=\sigma\boldsymbol E\)，则其导电区域 \(\boldsymbol E=0\Rightarrow\boldsymbol J=0\)。但超导体能在 \(\boldsymbol E=0\) 时有持久电流，非均匀电化学势和暂态也不同于简单欧姆模型。反之导体处于静电平衡时内部电场为零，而“无电流”并不自动意味着整个空间电势处处相等。

\paragraph{思考题 16.2}
同种金属电阻率 \(\rho_e\)、密度 \(\rho_m\)、长度 \(L\) 固定时 \(R=\rho_e L/S\)，质量 \(M=\rho_m SL\)。消去截面积 \(S=M/(\rho_m L)\)，有 \(\boxed{R=\rho_e\rho_m L^2/M\propto M^{-1}}\)。若温度、电阻率、材料纯度或长度不同，此结论不成立。

\paragraph{思考题 16.3}
本征半导体的热激发载流子数常近似 \(n_i\propto T^{3/2}\exp[-E_g/(2k_BT)]\)，导电率 \(\sigma=e(n\mu_e+p\mu_h)\)。虽迁移率 \(\mu\) 可随声子散射增多而下降，但 \(n,p\) 的指数增加可能主导，使 \(\boxed{R(T)\text{下降}}\)。对高温、杂质导电或相变条件，必须重新分析；不能推广为所有绝缘材料永远负温度系数。

\paragraph{思考题 16.4}
在静磁或准静态集中参数回路、忽略变化磁通的情况下，\(\boldsymbol E=-\nabla V\)，闭合回路 \(\oint\boldsymbol E\cdot d\boldsymbol l=0\)，于是电源电动势和各电压降代数和为零，即 Kirchhoff 第二定律 \(\sum \Delta V=0\)。若有变化磁通穿过回路，法拉第定律 \(\oint\boldsymbol E\cdot d\boldsymbol l=-d\Phi_B/dt\) 给出非零环流，须显式纳入感应电动势，不可无条件使用静电势回路和。

\paragraph{思考题 16.5}
电场是矢量 \(\boldsymbol E\)，定义为对试验电荷的力除以电荷；电势 \(V\) 是标量（静电下可定义），两者满足 \(\boxed{\boldsymbol E=-\nabla V}\)。零电势是参考点的选择；电场是否为零与电势\emph{梯度}有关。例如 \(V=0\) 的某等势面上仍可能有非零垂直电场；整个区域 \(V=\mathrm{const}\) 才有 \(\boldsymbol E=0\)。

\paragraph{思考题 16.6}
将未知高阻 \(R\) 与已知电容 \(C\) 串联，充至初始电压 \(V_0\) 后经过未知电阻放电，测量两时刻 \(V_1,V_2\) 与间隔 \(\Delta t\)。由 \(C\,dV/dt=-V/R\) 得 \(\boxed{R=\Delta t/[C\ln(V_1/V_2)]}\)。必须使用输入阻抗远高于待测 \(R\) 的电压测量仪、补偿漏电和寄生电容，否则时间常数被改变。

\paragraph{思考题 16.7}
原图16.18测量仪接在不同支路上，若电压表理想开路、电流表理想零内阻，其所串联的10 kΩ支路仍连通，并与其它电阻共享节点；端口AB的等效电阻可按图中真实连线求得，而不是仅由“表有高内阻”猜测。若给出\emph{非理想}仪表电阻则等效值会变化；因此必须先确认图中A、B之外电路所有连接点以及仪表约定。若按图理想V开路、A短路且中间唯一10 kΩ支路、上方5 kΩ串联，则 \(\boxed{R_{AB}=15\,\mathrm{k\Omega}}\)，这是假设具体理想仪表的结果。

\paragraph{思考题 16.8}
欧姆近似 \(I=U/R\)，若手到脚电阻 \(R\sim10\,\mathrm{k\Omega}\)，仅100V就可能产生 \(I\sim10\,\mathrm{mA}\)，危险程度受持续时间、路径、皮肤湿度和频率影响。故应使用绝缘工具、断电和检电，不可凭“电压不高”触摸电路。人体触电危害由穿过组织的电流、持续时间及路径共同决定，且人体阻抗\emph{不是恒定10 kΩ}。

\paragraph{思考题 16.9}
电源或尖端电离装置在传送带上喷射电荷，皮带运动将电荷搬运至高压金属球内壁，由梳状电极将电荷转移至球表面。克服由已积电荷形成的电势差而搬运新电荷需由电机等提供\emph{机械功}，故有效非静电力来自皮带机械输运及电晕充电源。单位电荷所作非静电功构成有效电动势 \(\mathcal E=W_{\rm nc}/q\)；能量守恒为外部机械功或外部电源功转成 \(\tfrac12 CV^2\) 等静电能，非静电力不是保守的库仑力。
